% Part: methods
% Chapter: proofs
% Section: using-definitions

\documentclass[../../../include/open-logic-section]{subfiles}

\begin{document}

\olfileid{mth}{prf}{def}

\olsection{استخدام التعريفات}

ذكرنا أنه يجب أن تكون ملمًا بكل التعريفات التي قد تُستخدم في البرهان،
وأن تستطيع تطبيقها تطبيقًا صحيحًا. وهذه نقطة مهمة حقًا، وجديرة بشيء
من التفصيل. تُستخدم التعريفات لاختصار الخواص والعلاقات كي نتحدث عنها
بإيجاز أكبر. ويسمى الاختصار المدخل \emph{المعرَّف}، وما يختصره
\emph{المُعرِّف}. وفي البراهين، كثيرًا ما نضطر إلى الرجوع إلى كيفية
إدخال المعرَّف، لأن علينا استغلال البنية المنطقية للمُعرِّف (أي النسخة
الطويلة التي يختصرها الحد المعرَّف) كي ننجز برهاننا. وبفك التعريفات،
تضمن بلوغك لبّ الموضع الذي يجري فيه العمل المنطقي.

سنبدأ بمثال. افترض أنك تريد إثبات ما يأتي:

\begin{prop}
لأي مجموعتين~$A$ و$B$، يكون~$A \cup B = B \cup A$.
\end{prop}

لكي نبدأ البرهان أصلًا، يجب أن نعرف معنى تطابق مجموعتين؛ أي يجب أن
نعرف معنى «$=$» في تلك المعادلة بالنسبة إلى المجموعات. وتُعرّف
المجموعات بأنها متطابقة متى كانت لها !!{element}s نفسها. ومن ثم يكون
التعريف الذي علينا فكه هو:

\begin{defn}
تكون المجموعتان~$A$ و$B$ \emph{متطابقتين}،~$A = B$، إذا وفقط إذا كان
كل !!{element} من~$A$ !!a{element} من~$B$، والعكس صحيح.
\end{defn}

يستخدم هذا التعريف~$A$ و$B$ نائبين عن مجموعتين اعتباطيتين. وما
يعرّفه---أي \emph{المعرَّف}---هو التعبير «$A = B$»، وذلك بإعطاء الشرط
الذي يصدق تحته~$A = B$. وهذا الشرط---«كل !!{element} من~$A$ هو
!!a{element} من~$B$، والعكس صحيح»---هو \emph{المُعرِّف}.\footnote{في
هذه الحالة خاصة---وهو أمر محير جدًا!---عندما~$A = B$، تكون المجموعتان
$A$ و$B$ مجموعة واحدة بعينها، مع أننا نستخدم حرفين مختلفين لها في
الطرف الأيسر والطرف الأيمن. لكن الطريقتين اللتين تتعين بهما تلك
المجموعة قد تختلفان، وهذا ما يجعل التعريف غير تافه.} ويقرر التعريف أن
$A = B$ تصدق إذا، وفقط إذا (ونختصر هذا بعبارة «إذا وفقط إذا»)، تحقق
الشرط.

عند تطبيق التعريف، يجب أن تطابق~$A$ و$B$ الواردتين فيه مع الحالة التي
تعالجها. وفي حالتنا، يعني هذا أنه كي تصدق~$A \cup B = B \cup A$، يجب
أن يكون كل~$z \in A \cup B$ في~$B \cup A$ أيضًا، والعكس صحيح. يؤدي
التعبير~$A \cup B$ في القضية دور~$A$ في التعريف، ويؤدي~$B \cup A$
دور~$B$. ولأن~$A$ و$B$ مستخدمتان في التعريف وفي عبارة القضية التي
نثبتها، لكن في استعمالين مختلفين، يجب أن تحذر كيلا تخلط بينهما.
فمثلًا، سيكون من الخطأ أن تظن أن بوسعك إثبات القضية ببيان أن كل
!!{element} من~$A$ هو !!a{element} من~$B$، والعكس صحيح---إذ يثبت ذلك
$A = B$، لا~$A \cup B = B \cup A$. (وفوق ذلك، لما كان يجوز أن تكون
$A$ و$B$ أي مجموعتين، فلن تقطع شوطًا بعيدًا، إذ قد تكونان مختلفتين
تمامًا إذا لم يُفترض شيء عن~$A$ و$B$.)

نتعامل داخل البرهان مع مفاهيم مجموعاتية مثل الاتحاد، ولذلك يجب أن
نعرف أيضًا معنى الرمز~$\cup$ كي نفهم كيفية متابعة البرهان. وأحيانًا
يولد فك تعريف ما تعريفات أخرى يلزم فكها. فمثلًا، تُعرّف
$A \cup B$ بأنها~$\Setabs{z}{z \in A \text{ أو } z \in B}$. فإذا
أردت إثبات~$x \in A \cup B$، أخبرك فك تعريف~$\cup$ بأن عليك إثبات
$x \in \Setabs{z}{z \in A \text{ أو } z \in B}$. وعليك الآن أيضًا
أن تتذكر أن~$x \in \Setabs{z}{\dots z\dots}$ إذا وفقط إذا كانت
$\dots x\dots$. ومن ثم، بمزيد من فك تعريف ترميز
$\Setabs{z}{\dots z \dots}$، يكون ما عليك بيانه هو:~$x \in A$ أو
$x \in B$. وهكذا، فإن «كل !!{element} من~$A \cup B$ هو أيضًا
!!a{element} من~$B \cup A$» تعني في الحقيقة: «لكل~$x$، إذا كان
$x \in A$ أو~$x \in B$، فإن~$x \in B$ أو~$x \in A$». وإذا فككنا
تعريفات القضية بالكامل، رأينا أن ما علينا بيانه هو:

\begin{prop}
لأي مجموعتين~$A$ و$B$: (أ) لكل~$x$، إذا كان~$x \in A$ أو
$x \in B$، فإن~$x \in B$ أو~$x \in A$؛ و(ب) لكل~$x$، إذا كان
$x \in B$ أو~$x \in A$، فإن~$x \in A$ أو~$x \in B$.
\end{prop}

والمهم أن فك التعريفات جزء ضروري من بناء البرهان. وقد يصعب إجراؤه
على الوجه الصحيح أحيانًا: إذ يجب أن تحرص على تمييز المتغيرات في
التعريف ومطابقتها مع الحدود في الادعاء الذي تثبته. ولكي تنجح، يجب أن
تعرف ما يطلبه السؤال ومعنى جميع الحدود المستخدمة فيه---وكثيرًا ما
تحتاج إلى فك أكثر من تعريف واحد. وفي البراهين البسيطة كالتي تأتي
أدناه، ينتج الحل مباشرة تقريبًا من التعريفات نفسها. لكنه لن يكون
دائمًا بهذه البساطة بالطبع.

\begin{prob}
افترض أنه طُلب منك إثبات~$A \cap B \neq \emptyset$. فك جميع التعريفات
الواردة هنا؛ أي أعد صياغة ذلك على نحو لا يذكر «$\cap$»، أو «$=$»، أو
«$\emptyset$».
\end{prob}

\end{document}
